\chapter{Introduction}
\label{ch:introduction}

\section{Background}

Consumer unmanned aerial vehicles (UAVs) are inexpensive, widely available, and increasingly implicated in unauthorised surveillance, perimeter incursion, and disruption of civil airspace operations. Counter-UAS (C-UAS) sensing modalities each have well-documented limitations~\cite{shi2018counteruas}: radar suffers on the small cross-section of quadcopter-class drones; radio-frequency detection requires an active command link and is defeated by autonomous pre-programmed flight; acoustic arrays degrade under wind, traffic, and aircraft noise. Visual detection with deep neural detectors is passive, low-cost, and context-rich, and is the modality adopted in this work.

The YOLO family of single-stage detectors~\cite{redmon2016yolo} as packaged in the Ultralytics framework~\cite{ultralytics2024} is the de-facto real-time architecture for drone detection. At sufficient input resolution, a YOLO detector trained on a competent drone corpus produces high drone-class precision \emph{and} high drone-class recall on standard benchmarks~\cite{svanstrom2021real,jiang2021antiuav}: the baseline RGB detector --- the Stage-1 reference against which the production confuser-injected detector (\texttt{ft4}, \S\ref{sec:rgb_detector}) is compared --- reaches drone $P=0.940$ and $R=0.961$ on Svanstr\"om at \texttt{imgsz=1280} (Table~\ref{tab:rgb_comparison}).
% [source: ledger=retrainedv2-recall-collapse; eval=rgb_svan_baseline; run=eval/diagnose_failures.py; cache=eval/results/_failure_diagnosis/svanstrom_1280_by_category.csv; config=svan_iop_1280]

Benchmark precision and recall are not deployment readiness. The same detector that scores $P=0.940$, $R=0.961$ on a curated split fails in the field for three reasons load-bearing in this thesis. The first is \emph{out-of-distribution confuser fire}: on the same dataset the baseline fires on 94.4\% of bird-only frames, 74.6\% of airplane-only frames, and 66.2\% of helicopter-only frames --- frames with no drone present, where every detection is by construction a false positive. The second is \emph{resolution}: small or distant drones fall below the resolvable floor at standard input sizes and disappear from the detector. The third is \emph{modality availability}: a single visible-light camera is fragile to lighting, glare, and clutter, and the thermal sensor that complements it is not always present in the field. This thesis builds a dual-modality detection system that closes these gaps and is honest about how each is measured. Confuser hallucination is the hardest of the three and accordingly takes the most space; but the contribution is the \emph{system}, and the discipline that builds it, not any single filter.
% [source: ledger=retrainedv2-recall-collapse; eval=rgb_svan_baseline; run=eval/diagnose_failures.py; cache=eval/results/_failure_diagnosis/svanstrom_1280_by_category.csv; config=svan_iop_1280]

\section{Problem Statement}

The simplest mitigation (retraining the detector with confuser images as hard negatives) has diminishing returns. The \texttt{hardneg\_v3more} variant of the baseline, trained with the Svanstr\"om bird/airplane/helicopter splits supplied as negatives, reduces helicopter hallucination from 66.2\% to 41.9\% but leaves bird hallucination essentially unchanged at 94.2\%, with airplane hallucination only dropping to 64.7\%. A more aggressive retraining (\texttt{retrained\_v2}) suppresses bird hallucination to 3.4\% but collapses drone recall on Svanstr\"om from $R=0.961$ to $R=0.306$ at the same $P=0.943$.
% [source: ledger=retrainedv2-recall-collapse; eval=rgb_svan_hardneg; run=eval/diagnose_failures_all.py; cache=eval/results/_failure_diagnosis/svanstrom_1280_by_category.csv; config=svan_iop_1280]

The Svanstr\"om collapse is dataset-specific. On the Roboflow OOD drone set the same \texttt{retrained\_v2} model retains $R=0.726$ versus the baseline's $R=0.746$, so the collapse reflects a small-drone resolution failure under aggressive confuser augmentation rather than a general defect of the model. For any deployment that must catch small or distant drones, however, the Svanstr\"om number is disqualifying: detector recall, once lost, cannot be recovered downstream --- no later stage adds detections the detector never proposed.
% [source: ledger=selcom-ood-confuser-damage; eval=rob_rgb_retrainedv2; run=eval/run_roboflow_eval.py; config=roboflow_rgb_drone_640]

Any discrimination the detector cannot learn without costing recall must be moved off it. This thesis therefore relocates confuser discrimination, modality arbitration, and false-positive filtering into a downstream dual-modality system, leaving the base detector tuned purely for recall.

\section{Research Questions}
\label{sec:rqs}

The thesis is organised around three research questions; answers are mapped to evidence in \S\ref{sec:rq_answers}.

\begin{description}
    \item[RQ1] Can a learned per-frame trust router, paired with a confuser filter that reuses the detector's own intermediate features, suppress out-of-distribution confuser fire rate substantially below what hard-negative mining alone achieves at the detector stage, without retraining the base detector?
    \item[RQ2] What is the in-distribution drone-detection cost of inserting that system, and how does it depend on the evaluation surface (paired-frame benchmark vs.\ real-world video) and on the choice of trust classifier?
    \item[RQ3] To what extent do detection features learned exclusively on thermal-infrared imagery transfer, zero-shot, to grayscale-converted visible-light video, and on which scenes does this transfer hold?
\end{description}

A methodological thread runs through all three: a \emph{statistics-before-training} discipline --- a model-agnostic feature-space instrument (the Model MRI, \S\ref{sec:model_mri}) and a leakage-aware feature-selection statistic that decide what a detector's features can separate before any training run --- together with a human-in-the-loop review process that co-develops the IR detector with its training corpus.

\section{Contributions}

This thesis makes four research contributions, resting on a shared methodological discipline described below.

\begin{enumerate}
    \item A \emph{deployable dual-modality drone-detection system}. An RGB and a thermal-infrared YOLO detector are fused by an XGBoost trust classifier that makes a per-frame four-way modality decision (\texttt{reject} / \texttt{trust\_rgb} / \texttt{trust\_ir} / \texttt{trust\_both}); the trusted modality's detections are then filtered by a per-frame confuser verifier --- the feature-reuse MLP of contribution~4, which superseded an earlier MobileNetV3 patch verifier on every surface measured. The figures below are reported for the configurations in which they were measured. An early classifier sweep with the open-world fallback classifier \texttt{fusion\_no\_fn\_v1.1} reduces the held-out OOD confuser zoo fire rate from 52.1\% (RGB detector alone) to 0.8\%, and the \texttt{scene\_aware\_v3more\_32feat} classifier yields 52.1\% $\to$ 10.3\% on the same chain. On Svanstr\"om paired data this OOD-robustness gain costs in-distribution drone F1: S1 trust-aware $F1=0.950$ (drone $P=0.940$, $R=0.961$) falls to $F1=0.895$ at the sa32 operating point. On the real-video diagnostic, with temporal smoothing, the cascade instead \emph{gains} drone segment $F1$ over the Stage-1 RGB number ($0.760 \to 0.826$, baseline RGB) while still cutting confuser segment FPR threefold. The in-distribution cost is surface-dependent, and the thesis reports both sides rather than only the favourable one.
    % [source: ledger=cascade-confuser-collapse; eval=clfzoo_fnfn; run=eval/cumulative_halluc.py; cache=eval/results/_cumulative_halluc/confuser_fusion_no_fn_model_v1.1/summary.json; config=confuser_zoo_1280]
    % [source: ledger=mlp-beats-patch-both-modalities; run=eval/overnight_ablation.py]
    \item A \emph{human-in-the-loop IR co-development case study}, reported in precision, recall, and dataset state across six revisions of the IR detector on a fixed held-out test split. The trace includes the V5 regression rather than hiding it, and identifies the corpus-level failure mode (uncurated positive ingestion) that caused it.
    \item An \emph{unexpected cross-modal transfer result}: the IR detector, trained zero-shot from thermal data, retains usable drone detection on grayscale-converted RGB video ($F1=0.636$ on 1{,}234 drone-positive frames across 9 YouTube clips; tied with the RGB baseline at $F1=0.837$ vs $0.840$ on the hardest single bird-cluttered clip while cutting confuser FPPI by $3.2\times$). The grayscale-RGB operating mode exists because thermal cameras are not always available in deployment; its strong performance was not predicted by design. The result is presented as a measured emergent property of training on a single intensity-channel modality.
    % [source: ledger=ir-grayscale-fallback,no-pareto-dominance-video; eval=vid_drone_ir_gray,vid_drone_baseline; run=eval/eval_video_tests.py; config=video_drone_iop] FPPI 3.2x = baseline 0.512 / ir_gray 0.158
    \item A \emph{feature-reuse confuser verifier}. Rather than re-processing detection crops with a separate MobileNetV3, a lightweight MLP consumes the detector's own fused intermediate ROI features (\texttt{p3}+\texttt{p5}), which a linear-discriminant probe shows are already nearly linearly separable ($\approx 95\%$ single-threshold accuracy on the shipped 32{,}931-detection corpus). On the production FT4 RGB detector (\S\ref{sec:rgb_detector}) the RGB verifier (\texttt{mlp\_v5}) beats the patch verifier on the confuser-rich surfaces (Svanstr\"om $F1$ $0.768\to0.869$; SelCom $+1.5$~pp; confuser-only hallucination rate $10.7\%\to0.8\%$, a $13\times$ reduction) and ties it on the saturated Anti-UAV surface. Reusing ROI features the detector has already computed, it runs $46$--$72\times$ faster per detection at $1$--$4\%$ pipeline overhead. The same recipe, with grayscale-harvested confusers z-aligned into thermal feature space, yields the production IR verifier (\texttt{mlp\_v5\_ir\_aligned}), which beats the patch verifier on the IR surfaces as well (held-out CBAM $F1$ $0.699\to0.846$, FP $48\to15$, recall-safe); across both modalities this verifier dominates the patch verifier on every surface measured. The one carve-out is a structural recall ceiling on the photo-style \texttt{rgb\_dataset} held-out split ($-11$~pp $F1$ vs the patch verifier), traced to a train/test distribution mismatch that three independent fixes failed to close; photo-style RGB is therefore routed to the patch verifier as a fallback.
    % [source: ledger=v5-beats-patch; eval=v5_svan_mlp; run=eval/eval_v4_vs_patch.py; cache=eval/results/_v5_selcom_pure_1x8; config=svan_iop_1280_s9]
    % [source: ledger=v5-rgbds-ceiling; eval=v5_rgbds_mlp; run=eval/eval_v4_vs_patch.py; config=rgb_dataset_iou_640]
\end{enumerate}

These four rest on one methodological discipline: \emph{statistics before training}. The Model MRI (\S\ref{sec:model_mri}) measures what a detector's features can separate before any training run; the same instrument yields the leakage statistic that selects the trust classifier's features and the linear-separability result that motivates the feature-reuse verifier above. Evaluation is held to the same standard: detections are scored under an explicit intersection-over-prediction rule (\S\ref{sec:metrics}), a Svanstr\"om usage/leakage audit (\S\ref{sec:svanstrom_audit}) bounds every in-distribution number, and a scoring-rule sensitivity check is reported --- dual versus trust-aware aggregation of identical multi-modal detections swings Svanstr\"om $F1$ by 28 percentage points, so any cross-study comparison must disclose its scoring rule. A reproducibility layer --- per-run JSON manifests, weight-hash-tagged inference caches, and a claim registry --- ties every reported number to the command that produced it.

\section{Ethical Considerations and Dual-Use}
\label{sec:ethics}

Counter-UAS sensing systems are dual-use technology. The same passive visual-detection capability that flags a hostile drone over critical infrastructure can, in a different deployment, contribute to mass aerial surveillance or armed-response targeting. This thesis develops the detection component only; it neither integrates with effectors (jammers, kinetic intercept, take-down systems) nor classifies operator behaviour. The training corpora contain no personally identifying information, no identifiable third-party imagery beyond what is publicly available in the cited datasets and YouTube clips, and no acoustic or RF data that could enable operator identification. The deployment scope assumed throughout is protection of fixed sites (perimeter monitoring, airport approaches, industrial facilities) where a detected drone is reported to a human operator who decides on response. The reproducibility layer (manifests, hash-tagged caches) is intended in part to make any future adaptation auditable.

\section{Thesis Outline}
\label{sec:outline}

Chapter~\ref{ch:literature} reviews the literature on real-time object detection, drone-detection datasets and the confuser problem, thermal-infrared imaging, multi-modal fusion, detection scoring, and cross-modal transfer. Chapter~\ref{ch:methodology} establishes the empirical infrastructure --- datasets, IoP scoring rule, the scoring-rule and Svanstr\"om leakage audits, training recipes, the reproducibility layer, and the Model MRI, defined before they are used --- and then presents the system architecture and component designs (\S\ref{ch:architecture}): the four-stage cascade, the fail-open design principle, trust-aware fusion, the alert-gate cascade, the temporal smoother, and the inline ablation justifying each production choice. Chapter~\ref{ch:hitl} reports the empirical evaluation: it first documents the dataset curation process --- the label-reviewer GUI, the HITL co-evolution loop, and the IR detector's six-revision trace including the V5 regression --- and then (\S\ref{ch:experiments}) reports cumulative confuser suppression, the Roboflow OOD audit, the real-video six-mode diagnostic, grayscale cross-modal transfer, the full cascade on real video, SelCom CCTV evaluation, comparisons with prior work, and threats to validity. Chapter~\ref{ch:conclusion} answers the three research questions, summarises the production stack, and lists future-work directions. Appendix~\ref{app:datasets} provides per-clip dataset metadata; Appendix~\ref{app:glossary} a glossary of abbreviations used throughout.

% ======================================================================
% CHAPTER 2 — RELATED WORK
% ======================================================================
